% Preamble
% =============================================================================

% Class of the document.
\documentclass[12pt,a4paper]{article}
% article : short article.
% report  : mid-length report.
% book    : book or thesis redaction.

% Paragraph skip length (default to 0).
\setlength{\parskip}{1ex}

% Packages
% =============================================================================

% Encoding
% -----------------------------------------------------------------------------

% Babel.
\usepackage[french]{babel}
% FontEnc.
\usepackage[T1]{fontenc}
% InputEnc.
\usepackage[utf8]{inputenc}

% Define \escapeus command to escape underscores.
\makeatletter
\DeclareRobustCommand*{\escapeus}[1]{
    \begingroup\@activeus\scantokens{#1\endinput}\endgroup}
\begingroup\lccode`\~=`\_\relax
    \lowercase{\endgroup\def\@activeus{\catcode`\_=\active \let~\_}}
\makeatother

% Text
% -----------------------------------------------------------------------------

% Acronym.
\usepackage{acronym}
% CsQuote.
\usepackage[style=french,french=guillemets]{csquotes}
% Enumerate.
\usepackage{enumerate}
% HyperRef.
\usepackage[hyperfootnotes=false,hidelinks]{hyperref}
% URL.
\usepackage{url}

% Algorithms
% -----------------------------------------------------------------------------

% Algorithm2E.
\usepackage[french,onelanguage,linesnumbered,ruled,vlined,commentsnumbered]{algorithm2e}

% Source code
% -----------------------------------------------------------------------------

% Listings.
\usepackage{listings}
% Minted.
\usepackage{minted}
% Caption.
\usepackage{caption}
\newenvironment{code}{\captionsetup{type=listing}}{}

% Files
% -----------------------------------------------------------------------------

% FancyVRB.
\usepackage{fancyvrb}
% Redefine \VerbatimInput.
\RecustomVerbatimCommand{\VerbatimInput}{VerbatimInput}
{
    fontsize=\footnotesize,
    frame=lines,         % Top and bottom rule only.
    framesep=1.5em,      % Separation between frame and text.
    rulecolor=\color{red!50!green!50!blue!50!},
    labelposition=topline,
    commandchars=\|\(\), % Escape character and argument delimiters for commands within the verbatim.
    commentchar=*        % Comment character.
}

% Figures
% -----------------------------------------------------------------------------

% GraphicX.
\usepackage{graphicx}
% SVG.
\usepackage{svg}
% WrapFig.
\usepackage{wrapfig}

% Charts
% -----------------------------------------------------------------------------

% PGFPLots
\usepackage{pgfplots}
\pgfplotsset{compat=1.16}
\usepgfplotslibrary{units}

% Mathematics
% -----------------------------------------------------------------------------

% AmsFonts.
\usepackage{amsfonts}
% AmsMath.
\usepackage{amsmath}
% AmsText.
\usepackage{amstext}
% AmsThm.
\usepackage{amsthm}
\newtheorem{prr}{Propriété}
\newtheorem{pro}{Proposition}
\newtheorem{thm}{Théorème}
\newtheorem{lem}{Lemme}
% NumPrint.
\usepackage{numprint}

% Physics
% -----------------------------------------------------------------------------

% Physics.
\usepackage{physics}

% Presentation
% -----------------------------------------------------------------------------

% XColor.
\usepackage{xcolor}

% References
% -----------------------------------------------------------------------------

% CleveRef.
\usepackage{cleveref}

% Structure.
% -----------------------------------------------------------------------------

% Geometry.
\usepackage{geometry}
% PDFLScape.
\usepackage{pdflscape}
% MultiCol.
\usepackage{multicol}
% TitleSec.
\usepackage{titlesec}
\newcommand{\sectionbreak}{\clearpage} % Use a page break before new sections.
% VMargin.
\usepackage{vmargin}
% FootMisc.
\usepackage[bottom]{footmisc}

% Symbols
% -----------------------------------------------------------------------------

% SIUnitX.
\usepackage{siunitx}

% Table
% -----------------------------------------------------------------------------

% Array.
\usepackage{array}
% BookTabs.
\usepackage{booktabs}
% CSVSimple.
\usepackage{csvsimple}

% Document
% =============================================================================

\begin{document}

\title{Analyse de performance et optimisation de code}
\author{AYOUB Pierre -- BONNAFOUS Camille -- FLAMANT Océane}

\maketitle

\begin{figure}[b]
    \centering
    \includegraphics[scale=0.3]{figures/isty.jpg}
\end{figure}

\newpage
\begin{abstract}

La simulation numérique est un procédé informatique visant à modéliser un
phénomène par ordinateur, s'agissant le plus souvent d'un phénomène
physique. Cette modélisation prend forme par des systèmes d'équations
décrivant l'état du système physique représenté à chaque instant. De
nombreux domaines scientifiques convergent vers la simulation
informatique, tel que certaines branches de la physique, de l'analyse
et de l'optimisation mathématique, ou encore le calcul haute
performance en informatique. Enfin, la simulation trouve naturellement
de nombreuses applications concernant des sujets variés, tel que la
simulation du climat et des évènements météorologiques, la simulation
d'essais nucléaires, de l'effet d'un médicament sur un corps, ou encore
des astres et de l'univers. Ce rapport s'articulera donc autour de l'analyse et
de l'optimisation d'un code de calcul, coeur des simulations numériques
présentés ci-dessus.

\end{abstract}

\tableofcontents

\section{Introduction}

Le projet que nous vous présentons aujourd'hui consiste à analyser puis,
grâce à nos mesures, optimiser un code de calcul, appelé kernel.

TODO

Le déroulement du projet s'est effectué en plusieurs étapes distinctes :
\begin{description}
    \item[Analyse du code] Cette phase consiste à analyser le programme d'un
        point de vue d'architecture informatique. Il convient d'étudier les
        choix mis en œuvres afin d'implémenter le ou les calculs nécessaires.
    \item[Protocole expérimental] Une fois l'analyse effectuée, nous
        pouvons en déduire le moyen le plus adapté afin de mesurer les
        performances de notre implémentation. Nous allons donc mettre en
        avant les critères théoriques à atteindre dans nos mesures, puis
        nous exposerons la manière dont nous avons mis ceci en pratique.
    \item[Optimisations et mesures] Grâce au protocole mis en place, nous
        pouvons quantifier la performance du programme. De ce fait, nous serons
        en mesure d'expérimenter différentes techniques d'optimisation sur le
        programme et d'en calculer l'accélération.
\end{description}

\section{Analyse du code}

\section{Protocole expérimental}

La mise en place d'un protocole expérimental de mesure est une étape nécessaire
et cruciale dans tout optimisation de code. D'une part, le but de ce protocole
est de mettre en lumière les points chauds du programme, c'est-à-dire les
parties du code qui ralentissent considérablement l'exécution de la simulation.
Ces points chauds seront les cibles de nos optimisations. D'autre part, après
chaque tentative d'optimisation, le protocole doit nous permettre de mesurer
l'impact de cette dernière, qu'il soit positif ou négatif, et enfin de le
quantifier.

TODO

\subsection{Théorie}

Lors de nos expériences, il ne faut pas oublier que le hasard ou l'aléa des
mesures peuvent biaiser un résultat. Afin d'éviter cela, il faut donc utiliser
une valeur moyenne ou une valeur médiane.

TODO

\subsection{Pratique}

TODO

\section{Optimisations et mesures}

\subsection{Phase 1}

\subsection{Phase 2}

\section{Conclusion}

TODO

\newpage
\section*{Acronymes}

\begin{acronym}
\end{acronym}

\end{document}
